% บทคัดย่อภาษาไทย
โครงงานฉบับนี้มีวัตถุประสงค์เพื่อออกแบบและพัฒนาเกมเดินเรือด้วยระบบสั่งการด้วยเสียง (Voice-Controlled Ship Navigation Game) สำหรับเป็นสื่อการเรียนรู้เพื่อส่งเสริมทักษะการสื่อสารและลดความวิตกกังวลในการใช้ภาษาต่างประเทศของผู้เรียน โดยระบบผสมผสานรูปแบบเกมจำลองสถานการณ์เข้ากับเทคโนโลยีการรู้จำเสียง (Speech Recognition) ประกอบด้วย 5 ระบบหลัก ได้แก่ ระบบรับคำสั่งเสียงและการประมวลผลคำสั่ง ระบบฟิสิกส์และการควบคุมยานพาหนะ ระบบลูกเรือปัญญาประดิษฐ์ ระบบสุ่มสร้างสิ่งกีดขวาง และระบบอินเทอร์เฟซผู้ใช้ รองรับผู้ใช้ในบทบาทกัปตันเรือ (ผู้เล่น) ที่ต้องสื่อสารและสั่งการลูกเรือปัญญาประดิษฐ์ ได้แก่ พลต้นหนและพลถือท้าย ระบบพัฒนาด้วยเอนจิน Unity 3D และภาษา C\# สำหรับโครงสร้างเกม และใช้เทคโนโลยีประมวลผลภาษาธรรมชาติร่วมกับโมเดลการรู้จำเสียงอย่าง Whisper AI สำหรับแปลงคำสั่งเสียงภาษาอังกฤษเป็นข้อความเพื่อประมวลผล ระบบลูกเรือปัญญาประดิษฐ์ช่วยให้เกิดการจำลองสถานการณ์สื่อสารขัดข้องเมื่อผู้เล่นออกเสียงไม่ชัดเจนหรือใช้คำสั่งผิดบริบท โดยลูกเรือจะทวนถามหรือขอให้ย้ำคำสั่งได้โดยตรง ทำให้การฝึกความคล่องแคล่วในการพูดและการแก้ปัญหาเฉพาะหน้าภายใต้ความกดดันมีประสิทธิภาพยิ่งขึ้น โดยจะมีการวัดผลการทดสอบและประเมินประสิทธิภาพผ่านแบบประเมินจากกลุ่มผู้ทดสอบจำนวน 40-50 คน